\chapter{REFLECTION AND CONCLUSION}

The development of the modular ZK-rollup benchmarking framework provided me with substantial technical experience across systems programming, blockchain settlement logic, benchmark automation, and distributed metric analysis. My work required me to understand not only how individual components function, but also how those components interact across the complete rollup pipeline. In particular, I gained practical experience in Rust-based service development, Layer-1 interaction using Ethereum client libraries, smart contract integration, data availability configuration, Docker-based benchmark orchestration, synthetic workload generation, and Python-based analysis of distributed system metrics.

A major technical learning outcome was the importance of clean abstraction boundaries in systems that interact with external infrastructure. The Submitter service originally had to coordinate several concerns at once: data availability formatting, Layer-1 transaction submission, receipt confirmation, error handling, and metrics collection. Through the introduction of abstractions such as the \texttt{BridgeClient} trait, I learned how decoupling infrastructure-specific logic from application-level orchestration improves testability and maintainability. This design allowed Submitter behaviour to be tested using mock bridge clients instead of requiring a live Layer-1 node for every test case.

I also developed a stronger understanding of reliability issues in blockchain-integrated systems. One important example was the Submitter receipt-status fix, where a transaction being mined was not sufficient evidence of successful execution. By adding explicit receipt-status checking for reverted Layer-1 transactions, the Submitter became more accurate in distinguishing successful settlement from failed contract execution. This was important for the validity of the benchmark results because a reverted transaction incorrectly recorded as successful would distort finalization and submission metrics.

Beyond the Submitter, I gained practical experience in designing empirical evaluation infrastructure. The benchmark suite required more than simply running services manually. It needed repeatable configuration, service startup ordering, workload generation, metrics collection, and output organization. Implementing scripts such as \path{run_experiment.sh}, \path{run_matrix.sh}, and \path{wait_for_sequencer.sh} showed me how small automation details, such as waiting for service readiness, can significantly affect the reliability of experimental results. Similarly, the Poisson-based workload generator and data-tools pipeline helped connect controlled input generation with measurable output analysis.

My work also highlighted the importance of being honest about contribution boundaries in a collaborative project. Some modules, especially the Executor and parts of the Sequencer, were developed through shared or mixed commits. Therefore, this report limited my claims to evidence-backed contributions such as Submitter implementation, L1 contract integration, benchmark orchestration, workload generation, data analysis tooling, and specific Sequencer robustness fixes. This evidence-based approach helped distinguish individual implementation work from overall team outcomes.

\vspace{1cm}
\noindent\textbf{Final Contribution Statement:}

During this final year project, my primary contribution was the implementation and integration of the Layer-1 interaction path and empirical evaluation infrastructure for the RollupX modular ZK-rollup benchmarking framework. I contributed to the Submitter service, including its application-level orchestration, data availability strategy handling, Layer-1 bridge interaction, receipt confirmation logic, metrics support, and mock-testable architecture through the \texttt{BridgeClient} abstraction. I also implemented reliability improvements such as explicit Layer-1 receipt-status checking to prevent reverted transactions from being recorded as successful submissions.

In addition to the Submitter, I contributed to the Layer-1 smart contract components used by the local benchmark environment, including bridge, data availability, and verifier-related contracts. These contracts provided the settlement boundary required for evaluating calldata, EIP-4844-style blob, and off-chain data availability configurations. I also developed the benchmark suite infrastructure, including experiment configuration, run orchestration scripts, service readiness handling, and repeatable execution support. The synthetic workload generator I contributed used configurable Poisson-style arrival timing to generate controlled transaction traffic for the Sequencer.

Finally, I built and integrated the Python-based data-tools pipeline used to aggregate distributed benchmark outputs, compute metrics such as throughput, latency percentiles, cost-related summaries, and data availability footprint, and generate analysis plots for the final evaluation. Together, these contributions helped connect separate project components into an end-to-end measurable research pipeline, from workload generation and batch processing to Layer-1 submission and post-experiment analysis. Through this work, I contributed both implementation functionality and experimental reliability to the RollupX benchmarking framework.

